\chapter{2005年福建卷（文）}
\section{单选题}
\begin{problem}
已知集合 \(P = \{x \mid |x - 1| \leqslant 1, x \in \R\}\)，集合 \(Q = \{x \mid x \in \N\}\)，则 \(P \cap Q\) 等于（\quad）

\choices
    {\(P\)}
    {\(Q\)}
    {\(\{1,2\}\)}
    {\(\{0,1,2\}\)}
\end{problem}

\begin{answer}
D
\end{answer}

\begin{solution}
由 \(\lvert x-1\rvert\leqslant1\)，得
\(-1\leqslant x-1\leqslant1\)，所以 \(0\leqslant x\leqslant2\)，即
\(P=[0,2]\). 又 \(Q=\N\)，而本题约定 \(0\in\N\)，因此 \(P\) 中的自然数只有
\(0\)，\(1\)，\(2\)，从而 \(P\cap Q=\{0,1,2\}\). 故选 D.
\end{solution}

\begin{problem}
不等式 \(\frac{2x - 1}{3x + 1} > 0\) 的解集是（\quad）

\choices
    {\(\left\{x\mid x < - \frac{1}{3}\text{\ 或\ } x > \frac{1}{2}\right\}\)}
    {\(\left\{x\mid -\frac{1}{3} < x < \frac{1}{2}\right\}\)}
    {\(\left\{x\mid x > \frac{3}{2}\right\}\)}
    {\(\left\{x\mid x > -\frac{1}{3}\right\}\)}
\end{problem}

\begin{answer}
A
\end{answer}

\begin{solution}
分式有意义的条件是 \(3x+1\neq0\)，即 \(x\neq-\frac{1}{3}\). 分子
\(2x-1\) 的零点为 \(x=\frac{1}{2}\)，所以两个临界点按从小到大排列为
\(-\frac{1}{3}\)，\(\frac{1}{2}\). 在区间
\(\left(-\infty,-\frac{1}{3}\right)\)，\(\left(-\frac{1}{3},\frac{1}{2}\right)\)，
\(\left(\frac{1}{2},+\infty\right)\) 内分别取 \(x=-1\)，\(x=0\)，\(x=1\)，
分式的符号依次为正、负、正. 由于不等式是严格不等式，\(x=-\frac{1}{3}\)
和 \(x=\frac{1}{2}\) 均不能取，因此解集为
\(\left\{x\mid x< -\frac{1}{3}\text{ 或 }x>\frac{1}{2}\right\}\). 故选 A.
\end{solution}

\begin{problem}
已知等差数列 \(\{a_{n}\}\) 中，\(a_{7} + a_{9} = 16\)，\(a_{4} = 1\)，则 \(a_{12}\) 的值是（\quad）

\choices
    {\(15\)}
    {\(30\)}
    {\(31\)}
    {\(64\)}
\end{problem}

\begin{answer}
A
\end{answer}

\begin{solution}
设等差数列的公差为 \(d\). 等差数列中，等距离两项的和等于中间项的两倍，
所以 \(a_{7}+a_{9}=2a_{8}=16\)，从而 \(a_{8}=8\). 由通项关系
\(a_{8}-a_{4}=(8-4)d\)，得
\(d=\frac{8-1}{4}=\frac{7}{4}\). 因而
\(a_{12}=a_{4}+(12-4)d=1+8\times\frac{7}{4}=15\). 故选 A.
\end{solution}

\begin{problem}
函数 \( y = \cos 2x \) 在下列哪个区间上是减函数（\quad）

\choices
    {\(\left[-\frac{\pi}{4}, \frac{\pi}{4}\right]\)}
    {\(\left[\frac{\pi}{4}, \frac{3\pi}{4}\right]\)}
    {\(\left[0, \frac{\pi}{2}\right]\)}
    {\(\left[\frac{\pi}{2}, \pi\right]\)}
\end{problem}

\begin{answer}
C
\end{answer}

\begin{solution}
函数的导数为 \(y'=-2\sin2x\). 当 \(0<x<\frac{\pi}{2}\) 时，
\(0<2x<\pi\)，所以 \(\sin2x>0\)，从而 \(y'<0\). 由导数符号可知，
函数在 \(\left[0,\frac{\pi}{2}\right]\) 上是减函数.

再检验其余选项：在 \(\left(-\frac{\pi}{4},0\right)\) 内，\(\sin2x<0\)，
所以函数递增；在 \(\left(\frac{\pi}{4},\frac{3\pi}{4}\right)\) 内，
\(\sin2x\) 在 \(x=\frac{\pi}{2}\) 两侧符号相反，函数先减后增；在
\(\left(\frac{\pi}{2},\pi\right)\) 内，\(\sin2x<0\)，函数递增. 因此只有选项 C 正确.
\end{solution}

\begin{problem}
下列结论正确的是（\quad）

\choices
    {当 \(x > 0\) 且 \(x \neq 1\) 时，\(\lg x + \frac{1}{\lg x} \geqslant 2\)}
    {当 \(x > 0\) 时，\(\sqrt{x} + \frac{1}{\sqrt{x}} \geqslant 2\)}
    {当 \(x \geqslant 2\) 时，\(x + \frac{1}{x}\) 的最小值为 2}
    {当 \(0 < x \leqslant 2\) 时，\(x - \frac{1}{x}\) 无最大值}
\end{problem}

\begin{answer}
B
\end{answer}

\begin{solution}
对选项 B，令 \(t=\sqrt{x}\)，则 \(t>0\). 因为
\(t+\frac{1}{t}-2=\frac{(t-1)^2}{t}\geqslant0\)，所以
\(\sqrt{x}+\frac{1}{\sqrt{x}}=t+\frac{1}{t}\geqslant2\). 当且仅当
\(t=1\)，即 \(x=1\) 时取等号，故选项 B 正确.

选项 A 不恒成立：取 \(\lg x=-1\)，即 \(x=10^{-1}>0\)，则
\(\lg x+\frac{1}{\lg x}=-1-1=-2<2\). 对选项 C，令
\(h(x)=x+\frac{1}{x}\)，则在 \(x\geqslant2\) 时
\(h'(x)=1-\frac{1}{x^2}>0\)，所以最小值在 \(x=2\) 处取得，为
\(h(2)=\frac{5}{2}\)，不是 2. 对选项 D，令
\(r(x)=x-\frac{1}{x}\)，则 \(r'(x)=1+\frac{1}{x^2}>0\)，
所以在 \(0<x\leqslant2\) 上最大值在 \(x=2\) 处取得，为 \(r(2)=\frac{3}{2}\)，
并非无最大值.
\end{solution}

\begin{problem}
若函数 \( f(x) = a^{x - b} \) 的图象如图，其中 \(a\)，\(b\) 为常数，则下列结论正确的是（\quad）

\bitmapfigure[width=0.34\linewidth]{img/2005/fujian_liberal/q6_fig1.png}

\choices
    {\(a > 1\)，\(b < 0\)}
    {\(a > 1\)，\(b > 0\)}
    {\(0 < a < 1\)，\(b > 0\)}
    {\(0 < a < 1\)，\(b < 0\)}
\end{problem}

\begin{answer}
D
\end{answer}

\begin{solution}
指数函数 \(y=a^{x-b}\) 的定义要求 \(a>0\) 且 \(a\neq1\). 图象从左向右递减，
因此 \(0<a<1\). 又因为 \(f(b)=a^{b-b}=1\)，所以图象与直线 \(y=1\) 的
交点横坐标就是 \(b\). 由题图可见该交点位于 \(y\) 轴左侧，因此 \(b<0\). 故选 D.
\end{solution}

\begin{problem}
已知直线 \(m\)，\(n\) 与平面 \(\alpha\)，\(\beta\)，给出下列三个命题：

\begin{circlelist}
\item 若 \(m \parallel \alpha\)，\(n \parallel \alpha\)，则 \(m \parallel n\)；
\item 若 \(m \parallel \alpha\)，\(n \perp \alpha\)，则 \(n \perp m\)；
\item 若 \(m \perp \alpha\)，\(m \parallel \beta\)，则 \(\alpha \perp \beta\).
\end{circlelist}

其中真命题的个数是（\quad）

\choices
    {\(0\)}
    {\(1\)}
    {\(2\)}
    {\(3\)}
\end{problem}

\begin{answer}
C
\end{answer}

\begin{solution}
命题一不成立. 例如取平面 \(\alpha\) 内两条相交直线的平行线，所得两条直线都平行于
\(\alpha\)，但彼此相交，因此不能推出 \(m\parallel n\).

命题二成立. 因为 \(n\perp\alpha\)，所以 \(n\) 垂直于平面 \(\alpha\) 内的每一条直线.
由 \(m\parallel\alpha\)，可在平面 \(\alpha\) 内作直线 \(m'\)，使 \(m'\parallel m\).
于是 \(n\perp m'\). 两条平行直线与同一直线所成的角相等，所以 \(n\perp m\).

命题三成立. 由 \(m\perp\alpha\)，直线 \(m\) 的方向垂直于平面 \(\alpha\). 又由
\(m\parallel\beta\)，可在平面 \(\beta\) 内作直线 \(m'\parallel m\)，从而
\(m'\perp\alpha\). 平面 \(\beta\) 含有一条垂直于平面 \(\alpha\) 的直线，故
\(\alpha\perp\beta\). 因此真命题有 \(2\) 个，选 C.
\end{solution}

\begin{problem}
已知 \(p: a \neq 0\)，\(q: ab \neq 0\)，则 \(p\) 是 \(q\) 的（\quad）

\choices
    {充分不必要条件}
    {必要不充分条件}
    {充要条件}
    {既不充分也不必要条件}
\end{problem}

\begin{answer}
B
\end{answer}

\begin{solution}
命题 \(q:ab\neq0\) 等价于 \(a\neq0\) 且 \(b\neq0\). 因此由 \(q\) 一定能推出
\(p:a\neq0\)，所以 \(p\) 是 \(q\) 的必要条件. 但反过来不成立：取
\(a=1\)，\(b=0\)，则 \(p\) 成立而 \(ab=0\)，所以 \(q\) 不成立. 故 \(p\) 不是
\(q\) 的充分条件，选 B.
\end{solution}

\begin{problem}
已知定点 \(A\)，\(B\) 且 \(|AB| = 4\)，动点 \(P\) 满足 \(|PA| - |PB| = 3\)，则 \(|PA|\) 的最小值是（\quad）

\choices
    {\( \frac{1}{2} \)}
    {\( \frac{3}{2} \)}
    {\( \frac{7}{2} \)}
    {\(5\)}
\end{problem}

\begin{answer}
C
\end{answer}

\begin{solution}
设 \(|PA|=r\). 由 \(|PA|-|PB|=3\)，得 \(|PB|=r-3\)，所以必有
\(r\geqslant3\). 另一方面，由三角不等式
\(|PA|+|PB|\geqslant|AB|=4\)，代入上式得
\(r+(r-3)\geqslant4\)，即 \(2r\geqslant7\)，从而
\(r\geqslant\frac{7}{2}\).

当 \(P\) 取在线段 \(AB\) 上，并满足 \(|PA|=\frac{7}{2}\)，\(|PB|=\frac{1}{2}\) 时，
有 \(|PA|+|PB|=|AB|\) 且 \(|PA|-|PB|=3\)，所以题设条件满足. 下界可以取得，
故 \(|PA|\) 的最小值为 \(\frac{7}{2}\)，选 C.
\end{solution}

\begin{problem}
从 \(6\) 人中选 \(4\) 人分别到巴黎，伦敦，悉尼，莫斯科四个城市游览，要求每个城市有一人游览，每人只游览一个城市，且这 \(6\) 人中甲，乙两人不去巴黎游览，则不同的选择方案共有（\quad）

\choices
    {\(300\) 种}
    {\(240\) 种}
    {\(144\) 种}
    {\(96\) 种}
\end{problem}

\begin{answer}
B
\end{answer}

\begin{solution}
巴黎的游览者不能是甲或乙，只能从其余 \(4\) 人中选取，因此有 \(4\) 种选法.
选定巴黎游览者后，还剩 \(5\) 人. 从这 \(5\) 人中选出分别去伦敦、悉尼、莫斯科的
\(3\) 人并安排城市，等价于从 \(5\) 人中取 \(3\) 人排列，方法数为
\(\mathrm A_{5}^{3}=5\times4\times3=60\). 因而总方案数为
\(4\mathrm A_{5}^{3}=4\times60=240\)，故选 B.
\end{solution}

\begin{problem}
如图，长方体 \(ABCD - A_{1}B_{1}C_{1}D_{1}\) 中，\(AA_{1} = AB = 2\)，\(AD = 1\)，点 \(E\)，\(F\)，\(G\) 分别是 \(DD_{1}\)，\(AB\)，\(CC_{1}\) 的中点，则异面直线 \(A_{1}E\) 与 \(GF\) 所成的角是（\quad）

\bitmapfigure[width=0.42\linewidth]{img/2005/fujian_liberal/q11_fig1.png}

\choices
    {\(\arccos \frac{\sqrt{15}}{5}\)}
    {\(\frac{\pi}{4}\)}
    {\(\arccos \frac{\sqrt{10}}{5}\)}
    {\(\frac{\pi}{2}\)}
\end{problem}

\begin{answer}
D
\end{answer}

\begin{solution}
建立空间直角坐标系，令 \(A(0,0,0)\)，\(B(2,0,0)\)，\(D(0,1,0)\)，
\(A_{1}(0,0,2)\). 于是 \(C(2,1,0)\)，并由中点公式得到
\(E(0,1,1)\)，\(F(1,0,0)\)，\(G(2,1,1)\).
因此可取两异面直线的方向向量为
\(\overrightarrow{A_{1}E}=(0,1,-1)\) 和 \(\overrightarrow{GF}=(-1,-1,-1)\). 计算得
\(\overrightarrow{A_{1}E}\cdot\overrightarrow{GF}=0\times(-1)+1\times(-1)+(-1)\times(-1)=0\).
所以两直线的方向垂直，异面直线 \(A_{1}E\) 与 \(GF\) 所成的角为
\(\frac{\pi}{2}\)，选 D.
\end{solution}

\begin{problem}
\(f(x)\) 是定义在 \(\R\) 上的以 \(3\) 为周期的偶函数，且 \(f(2) = 0\)，则方程 \(f(x) = 0\) 在区间 \((0,6)\) 内解的个数的最小值是（\quad）

\choices
    {\(5\)}
    {\(4\)}
    {\(3\)}
    {\(2\)}
\end{problem}

\begin{answer}
B
\end{answer}

\begin{solution}
由 \(f(2)=0\) 及偶性，得 \(f(-2)=f(2)=0\). 因为周期为 \(3\)，所以
\(f(1)=f(-2)=0\)，再由周期性得 \(f(4)=f(1)=0\)，同时
\(f(5)=f(2)=0\). 这四个数互不相同且都在 \((0,6)\) 内，因此零点个数至少为
\(4\).

下面构造一个恰好只有这四个零点的函数. 令
\(f(x)=\cos\frac{2\pi x}{3}+\frac{1}{2}\). 该函数为偶函数，且
\(f(x+3)=f(x)\)，所以以 \(3\) 为周期；并且
\(f(2)=\cos\frac{4\pi}{3}+\frac{1}{2}=0\). 其零点满足
\(\cos\frac{2\pi x}{3}=-\frac{1}{2}\)，即
\(x=1+3k\) 或 \(x=2+3k\)，其中 \(k\in\Z\). 在 \((0,6)\) 内零点恰为
\(1\)，\(2\)，\(4\)，\(5\)，共 \(4\) 个. 因而最小值为 \(4\)，选 B.
\end{solution}

\section{填空题}
\begin{problem}
\(\left(2\sqrt{x} -\frac{1}{x}\right)^6\) 展开式中的常数项是\fillinblank{}.
\end{problem}

\begin{answer}
240
\end{answer}

\begin{solution}
在二项式定理中，取 \(\left(-\frac{1}{x}\right)^k\) 的项时，通项为
\(\mathrm C_{6}^{k}(2\sqrt{x})^{6-k}\left(-\frac{1}{x}\right)^{k} =\mathrm C_{6}^{k}2^{6-k}(-1)^k x^{\frac{6-k}{2}-k} =\mathrm C_{6}^{k}2^{6-k}(-1)^k x^{3-\frac{3k}{2}}\)，\(k=0\)，\(1\)，\(\ldots\)，\(6\).
常数项的幂次必须为零，因此 \(3-\frac{3k}{2}=0\)，解得 \(k=2\). 代入通项，
常数项为 \(\mathrm C_{6}^{2}2^{4}(-1)^2=15\times16=240\).
\end{solution}

\begin{problem}
在 \(\triangle ABC\) 中，\(\angle C = 90^{\circ}\)，\(\overrightarrow{AB} = (k,1)\)，\(\overrightarrow{AC} = (2,3)\)，则 \(k\) 的值是\fillinblank{}.

\end{problem}

\begin{answer}
5
\end{answer}

\begin{solution}
因为 \(\angle C=90^{\circ}\)，所以 \(\overrightarrow{CA}\perp\overrightarrow{CB}\)，即
\(\overrightarrow{CA}\cdot\overrightarrow{CB}=0\). 由向量加法，
\(\overrightarrow{CA}=-\overrightarrow{AC}=(-2,-3)\)，而
\(\overrightarrow{CB}=\overrightarrow{AB}-\overrightarrow{AC}=(k-2,-2)\). 因而
\[
\left(-2,-3\right)\cdot(k-2,-2)=-2(k-2)+6=-2k+10=0
\]
解得 \(k=5\).
\end{solution}

\begin{problem}
非负实数 \(x\)，\(y\) 满足 \(\begin{cases}
2x + y \leqslant 0,\\
x + y - 3 \leqslant 0,
\end{cases}\) 则 \(x + 3y\) 的最大值为\fillinblank{}.

\end{problem}

\begin{answer}
0
\end{answer}

\begin{solution}
因为 \(x\geqslant0\)，\(y\geqslant0\)，所以 \(2x+y\geqslant0\). 结合条件
\(2x+y\leqslant0\)，可知 \(2x+y=0\). 两项 \(2x\) 与 \(y\) 都是非负数，
它们的和为零只能分别为零，因此 \(x=0\)，\(y=0\). 代入第二个不等式得
\(0+0-3=-3\leqslant0\)，所以该点确实可行. 此时
\(x+3y=0\)，可行域只有这一点，故最大值为 \(0\).
\end{solution}

\begin{problem}
把下面不完整的命题补充完整，并使之成为真命题：若函数 \( f(x)=3+\log_{2}x \) 的图象与 \( g(x) \) 的图象关于\fillinblank{} 对称，则函数 \( g(x)=\fillinblank{} \). 注：填上你认为可以成为真命题的一件情形即可，不必考虑所有可能的情形.
\end{problem}

\begin{answer}
\(x\) 轴；\(g(x)=-3-\log_{2}x\)
\end{answer}

\begin{solution}
取对称轴为 \(x\) 轴. 图象上任意一点 \((x,f(x))\) 关于 \(x\) 轴的对称点为
\((x,-f(x))\)，所以对应函数满足 \(g(x)=-f(x)\). 由于
\(f(x)=3+\log_{2}x\) 的定义域为 \(x>0\)，故
\(g(x)=-(3+\log_{2}x)=-3-\log_{2}x\)，\(x>0\).
因此可填“\(x\) 轴”和 \(g(x)=-3-\log_{2}x\)，该命题成为真命题.
\end{solution}

\section{解答题}

\begin{problem}
已知 \(-\frac{\pi}{2} < x < 0\)，\(\sin x + \cos x = \frac{1}{5}\).

\begin{enumerate}
\item 求 \( \sin x - \cos x \) 的值；
\item 求 \( \frac{\sin 2x + 2\sin^{2}x}{1 - \tan x} \) 的值.
\end{enumerate}

\end{problem}

\begin{solution}
\begin{enumerate}
\item 记 \(s=\sin x\)，\(c=\cos x\). 因为 \(-\frac{\pi}{2}<x<0\)，所以
\(s<0\)，\(c>0\)，且 \(s^2+c^2=1\). 由 \(s+c=\frac{1}{5}\)，有
\[
(s+c)^2=s^2+c^2+2sc=\frac{1}{25}
\]
从而 \(2sc=\frac{1}{25}-1=-\frac{24}{25}\)，即
\(sc=-\frac{12}{25}\). 因此
\[
(s-c)^2=s^2+c^2-2sc=1+\frac{24}{25}=\frac{49}{25}
\]
又因 \(s<0<c\)，所以 \(s-c<0\)，故 \(s-c=-\frac{7}{5}\). 因而
\(\sin x-\cos x=-\frac{7}{5}\).
\item 先化简分子. 由 \(\sin2x=2sc\) 和 \(s+c=\frac{1}{5}\)，得
\[
\sin2x+2\sin^2x=2sc+2s^2=2s(s+c)=\frac{2s}{5}
\]
又因为 \(c>0\)，所以 \(1-\tan x\) 有意义，并且
\[
1-\tan x=1-\frac{s}{c}=\frac{c-s}{c}
=\frac{\frac{7}{5}}{c}=\frac{7}{5c}
\]
于是
\[
\frac{\sin2x+2\sin^2x}{1-\tan x}
=\frac{2s}{5}\cdot\frac{5c}{7}
=\frac{2sc}{7}
=\frac{2}{7}\left(-\frac{12}{25}\right)
=-\frac{24}{175}
\]
\end{enumerate}
\end{solution}

\begin{problem}
甲，乙两人在罚球线投球命中的概率分别为 \(\frac{1}{2}\) 与 \(\frac{2}{5}\)，投中得 1 分，投不中得 0 分.

\begin{enumerate}
\item 甲，乙两人在罚球线各投球一次，求两人得分之和 \(\xi\) 的数学期望；
\item 甲，乙两人在罚球线各投球二次，求这四次投球中至少一次命中的概率.
\end{enumerate}
\end{problem}

\begin{solution}
\begin{enumerate}
\item 设甲、乙各投一次所得分数分别为 \(X\)，\(Y\)，则 \(\xi=X+Y\). 甲投中和乙投中的概率分别为
\(\frac{1}{2}\) 和 \(\frac{2}{5}\)，所以
\[
P(\xi=0)=\frac{1}{2}\cdot\frac{3}{5}=\frac{3}{10}
\]
\[
P(\xi=2)=\frac{1}{2}\cdot\frac{2}{5}=\frac{1}{5}
\]
而恰有一人投中的事件由“甲中乙不中”和“甲不中乙中”组成，故
\[
P(\xi=1)=\frac{1}{2}\cdot\frac{3}{5}
+\frac{1}{2}\cdot\frac{2}{5}=\frac{1}{2}
\]
因此
\[
E(\xi)=0\cdot\frac{3}{10}+1\cdot\frac{1}{2}
+2\cdot\frac{1}{5}=\frac{9}{10}
\]
\item 各次投球相互独立. 甲两次均未命中的概率为
\(\left(1-\frac{1}{2}\right)^2=\frac{1}{4}\)，乙两次均未命中的概率为
\(\left(1-\frac{2}{5}\right)^2=\frac{9}{25}\). 因而四次投球均未命中的概率为
\[
\frac{1}{4}\cdot\frac{9}{25}=\frac{9}{100}
\]
所求事件是“至少一次命中”，与“四次均未命中”互为对立事件，所以所求概率为
\[
1-\frac{9}{100}=\frac{91}{100}
\]
\end{enumerate}
\end{solution}

\begin{problem}
已知 \(\{a_{n}\}\) 是公比为 \(q\) 的等比数列，且 \(a_{1}\)，\(a_{3}\)，\(a_{2}\) 成等差数列.

\begin{enumerate}
\item 求 \(q\) 的值；
\item 设 \(\{b_n\}\) 是以 2 为首项，\(q\) 为公差的等差数列，其前 \(n\) 项和为 \(S_n\)，当 \(n \geqslant 2\) 时，比较 \(S_n\) 与 \(b_n\) 的大小，并说明理由.
\end{enumerate}
\end{problem}

\begin{solution}
\begin{enumerate}
\item 由等比数列的通项关系，\(a_2=a_1q\)，\(a_3=a_1q^2\). 因为 \(a_1\)，\(a_3\)，\(a_2\)
成等差数列，且中间项是 \(a_3\)，所以
\[
2a_3=a_1+a_2
\]
代入 \(a_2=a_1q\)，\(a_3=a_1q^2\)，得
\(a_1(2q^2-q-1)=0\). 等比数列的首项不能为零，所以 \(a_1\ne0\)，从而
\[
2q^2-q-1=(2q+1)(q-1)=0
\]
所以 \(q=1\) 或 \(q=-\frac{1}{2}\).
\item 等差数列 \(\{b_n\}\) 的首项为 \(2\)，公差为 \(q\)，故
\(b_n=2+(n-1)q\)，且
\(S_n=\frac{n}{2}\left[2b_1+(n-1)q\right]
=\frac{n}{2}\left[4+(n-1)q\right]\).

当 \(q=1\) 时，\(b_n=n+1\)，\(S_n=\frac{n(n+3)}{2}\)，因此
\[
S_n-b_n=\frac{n(n+3)}{2}-(n+1)=\frac{(n-1)(n+2)}{2}>0
\]
对一切 \(n\geqslant2\) 成立，所以 \(S_n>b_n\).

当 \(q=-\frac{1}{2}\) 时，\(b_n=2-\frac{n-1}{2}=\frac{5-n}{2}\)，且
\(S_n=\frac{n(9-n)}{4}\). 因而
\[
S_n-b_n=\frac{n(9-n)}{4}-\frac{5-n}{2}
=\frac{(n-1)(10-n)}{4}
\]
当 \(n\geqslant2\) 时，\(n-1>0\)，所以差值的符号由 \(10-n\) 决定. 因此
当 \(2\leqslant n\leqslant9\) 时 \(S_n>b_n\)，当 \(n=10\) 时 \(S_n=b_n\)，
当 \(n\geqslant11\) 时 \(S_n<b_n\).
\end{enumerate}
\end{solution}

\begin{problem}
已知函数 \( f(x) = x^3 + bx^2 + ax + d \) 的图象过点 \( P(0,2) \)，且在点 \( M(-1, f(-1)) \) 处的切线方程为 \( 6x - y + 7 = 0 \).

\begin{enumerate}
\item 求函数 \( y = f(x) \) 的解析式；
\item 求函数 \( y = f(x) \) 的单调区间.
\end{enumerate}
\end{problem}

\begin{solution}
\begin{enumerate}
\item 因为图象过点 \(P(0,2)\)，所以 \(f(0)=d=2\). 切线方程
\(6x-y+7=0\) 可写成 \(y=6x+7\). 点 \(M\) 的横坐标为 \(-1\)，所以
\(f(-1)=6(-1)+7=1\)，且切线斜率给出 \(f'(-1)=6\). 另一方面，
\(f(-1)=-1+b-a+d\)，\(f'(x)=3x^2+2bx+a\)
所以
\(-1+b-a+2=1\)，\(3-2b+a=6\).
即 \(b-a=0\)，\(a-2b=3\). 由 \(a=b\) 代入第二式得 \(-b=3\)，
从而 \(a=b=-3\). 因此
\[
f(x)=x^3-3x^2-3x+2
\]
\item 对上式求导，得
\[
f'(x)=3x^2-6x-3=3\left[(x-1)^2-2\right]
\]
令 \(f'(x)=0\)，得 \(x=1-\sqrt{2}\) 或 \(x=1+\sqrt{2}\). 因为
\(f'(x)\) 是开口向上的二次函数，所以
\[
f'(x)>0\quad\text{当 }x<1-\sqrt{2}\text{ 或 }x>1+\sqrt{2}
\]
\[
f'(x)<0\quad\text{当 }1-\sqrt{2}<x<1+\sqrt{2}
\]
所以函数在 \((-\infty,1-\sqrt{2})\) 和 \((1+\sqrt{2},+\infty)\) 上单调递增，
在 \((1-\sqrt{2},1+\sqrt{2})\) 上单调递减.
\end{enumerate}
\end{solution}

\begin{problem}
如图，直二面角 \(D - AB - E\) 中，四边形 \(ABCD\) 是边长为 2 的正方形，\(AE = EB\)，\(F\) 为 \(CE\) 上的点，且 \(BF \perp\) 平面 \(ACE\).

\begin{enumerate}
\item 求证 \(AE \perp\) 平面 \(BCE\)；
\item 求二面角 \(B - AC - E\) 的大小；
\item 求点 \(D\) 到平面 \(ACE\) 的距离.
\end{enumerate}

\bitmapfigure[width=0.42\linewidth]{img/2005/fujian_liberal/q21_fig1.png}

\end{problem}

\begin{solution}
\begin{enumerate}
\item 取空间直角坐标系，使正方形 \(ABCD\) 所在平面为 \(y=0\)，且平面 \(ABE\)
所在的另一面为 \(z=0\). 令
\(A(0,0,0)\)，\(B(2,0,0)\)，\(C(2,0,2)\)，\(D(0,0,2)\).
由 \(AE=EB\) 可知 \(E\) 在 \(AB\) 的垂直平分面内，又因 \(E\in z=0\)，
故可设 \(E(1,t,0)\).

设 \(F=C+\lambda(E-C)\). 则
\(F=(2-\lambda,\lambda t,2-2\lambda)\)，\(\overrightarrow{BF}=(-\lambda,\lambda t,2-2\lambda)\).
又
\(\overrightarrow{AC}=(2,0,2)\)，\(\overrightarrow{AE}=(1,t,0)\). 由
\(BF\perp\) 平面 \(ACE\)，向量 \(\overrightarrow{BF}\) 必垂直于该平面内相交的两条直线
\(AC\)，\(AE\)，所以
\[
\begin{cases}
\overrightarrow{BF}\cdot\overrightarrow{AC}=4-6\lambda=0,\\
\lambda(t^2-1)=0
\end{cases}
\]
第一式给出 \(\lambda=\frac{2}{3}\)，因而 \(\lambda\ne0\)，第二式给出
\(t^2=1\). 改变 \(y\) 轴方向可将 \(t=-1\) 化为 \(t=1\)，以下取
\(E(1,1,0)\). 此时
\(\overrightarrow{AE}=(1,1,0)\)，\(\overrightarrow{BE}=(-1,1,0)\)，\(\overrightarrow{CE}=(-1,1,-2)\).
直接计算得
\(\overrightarrow{AE}\cdot\overrightarrow{BE}=0\)，
\(\overrightarrow{AE}\cdot\overrightarrow{CE}=0\). 直线 \(BE\) 与 \(CE\) 在点
\(E\) 相交且都在平面 \(BCE\) 内，所以由直线垂直平面的判定定理，
\(AE\perp\) 平面 \(BCE\).
\item 平面 \(ABC\) 即 \(y=0\)，可取其法向量
\(\bs{n}_1=(0,1,0)\). 平面 \(ACE\) 的两条相交方向向量为
\(\overrightarrow{AC}=(2,0,2)\) 和 \(\overrightarrow{AE}=(1,1,0)\)，故可取法向量
\[
\bs{n}_2=\overrightarrow{AC}\times\overrightarrow{AE}=(-2,2,2)
\]
设二面角 \(B-AC-E\) 的平面角为 \(\theta\). 取两法向量的夹角的锐角，得
\[
\cos\theta
=\frac{|\bs{n}_1\cdot\bs{n}_2|}{|\bs{n}_1||\bs{n}_2|}
=\frac{2}{1\cdot2\sqrt{3}}
=\frac{\sqrt{3}}{3}
\]
所以二面角 \(B-AC-E\) 的大小为 \(\arccos\frac{\sqrt{3}}{3}\).
\item 由法向量 \((-2,2,2)\)，且平面 \(ACE\) 经过原点 \(A\)，可得平面方程
\(-2x+2y+2z=0\)，即 \(-x+y+z=0\). 因而点 \(D(0,0,2)\) 到该平面的距离为
\[
\frac{\lvert-0+0+2\rvert}{\sqrt{(-1)^2+1^2+1^2}}
=\frac{2}{\sqrt{3}}
=\frac{2\sqrt{3}}{3}
\]
\end{enumerate}
\end{solution}

\begin{problem}
已知方向向量为 \(\bs{v} = (1, \sqrt{3})\) 的直线 \(l\) 过点 \((0, -2\sqrt{3})\) 和椭圆 \(C: \frac{x^2}{a^2} + \frac{y^2}{b^2} = 1\)（\(a > b > 0\)）的焦点，且椭圆 \(C\) 的中心关于直线 \(l\) 的对称点在椭圆 \(C\) 的右准线上.

\begin{enumerate}
\item 求椭圆 \(C\) 的方程；
\item 是否存在过点 \(E(-2,0)\) 的直线 \(m\) 交椭圆 \(C\) 于点 \(M\)，\(N\)，满足 \(\overrightarrow{OM} \cdot \overrightarrow{ON} = \frac{4}{3}\sqrt{6}\cot \angle MON \neq 0\)（\(O\) 为原点）. 若存在，求直线 \(m\) 的方程；若不存在，请说明理由.
\end{enumerate}

\end{problem}

\begin{solution}
\begin{enumerate}
\item 由方向向量 \(\bs v=(1,\sqrt3)\) 和所过点 \((0,-2\sqrt3)\)，直线 \(l\) 的方程为
\[
y=\sqrt3x-2\sqrt3
\]
它与 \(x\) 轴交于 \((2,0)\). 椭圆的焦点在 \(x\) 轴上，且直线 \(l\) 经过一个焦点，
所以该焦点为 \((2,0)\)，从而 \(c=2\). 将直线改写为
\[
\sqrt3x-y-2\sqrt3=0
\]
设原点为 \(O\). 由点到直线的垂足公式，\(O\) 在该直线上的垂足为
\[
H=\left(\frac{3}{2},-\frac{\sqrt3}{2}\right)
\]
因此 \(O\) 关于 \(l\) 的对称点为
\[
O'=2H=\left(3,-\sqrt3\right)
\]
椭圆的右准线方程为 \(x=\frac{a^2}{c}=\frac{a^2}{2}\). 因为 \(O'\) 在右准线上，
\(a^2/2=3\)，所以 \(a^2=6\). 再由 \(c^2=a^2-b^2\)，得
\(b^2=6-4=2\). 故椭圆方程为
\[
\frac{x^2}{6}+\frac{y^2}{2}=1
\]
\item 记 \(K=\frac{4}{3}\sqrt6\)，\(\theta=\angle MON\)，并令
\(D=\overrightarrow{OM}\cdot\overrightarrow{ON}\)，\(\Delta=\left|\det\left(\overrightarrow{OM},\overrightarrow{ON}\right)\right|\).
因题设要求 \(D\neq0\)，且
\(\cot\theta=\frac{D}{\Delta}\)，原条件等价于
\(D=K\frac{D}{\Delta}\)，\(D\neq0\)
所以必须且只需满足 \(\Delta=K\).

先考虑非竖直直线. 设其斜率为 \(k\)，由于直线经过 \(E(-2,0)\)，其方程为
\(y=k(x+2)\). 令
\(x=-2+t\)，\(y=kt\).
代入椭圆方程，得
\[
(1+3k^2)t^2-4t-2=0
\]
设两个交点对应的根为 \(t_1\)，\(t_2\). 该二次方程的判别式为
\(16+8(1+3k^2)=24(1+k^2)>0\)，所以确有两个不同交点，且
\[
|t_1-t_2|=\frac{2\sqrt{6}\sqrt{1+k^2}}{1+3k^2}
\]
取方向向量 \(\bs u=(1,k)\)，则
\(M=E+t_1\bs u\)，\(N=E+t_2\bs u\). 因此
\[
\Delta=|t_1-t_2|\left|\det(E,\bs{u})\right|
=\frac{4\sqrt{6}|k|\sqrt{1+k^2}}{1+3k^2}
\]
令 \(\Delta=K\)，约去正数 \(4\sqrt6\)，得到
\[
3|k|\sqrt{1+k^2}=1+3k^2
\]
两边均非负，平方后有
\[
9k^2(1+k^2)=(1+3k^2)^2
\]
化简为 \(3k^2=1\)，所以
\(k=\pm\frac{\sqrt3}{3}\). 还需检验题设中的非零条件. 由韦达定理，
\(t_1+t_2=\frac{4}{1+3k^2}\)，\(t_1t_2=-\frac{2}{1+3k^2}\).
利用 \(E=(-2,0)\)，\(\bs u=(1,k)\)，可得
\[
D=\left(E+t_1\bs u\right)\cdot\left(E+t_2\bs u\right)
=4-2(t_1+t_2)+(1+k^2)t_1t_2
=\frac{10k^2-6}{1+3k^2}
\]
当 \(k^2=\frac13\) 时，\(D=-\frac43\neq0\)，所以两条非竖直直线均满足题设条件.

再考虑竖直直线 \(x=-2\). 代入椭圆方程，得到
\(y^2=\frac23\)，所以交点为
\(M\left(-2,\sqrt{\frac{2}{3}}\right)\)，\(N\left(-2,-\sqrt{\frac{2}{3}}\right)\).
此时
\[
D=\left(-2,\sqrt{\frac23}\right)\cdot
\left(-2,-\sqrt{\frac23}\right)
=4-\frac23=\frac{10}{3}\neq0
\]
且
\[
\Delta=\left|\det\left(
\left(-2,\sqrt{\frac23}\right),
\left(-2,-\sqrt{\frac23}\right)\right)\right|
=4\sqrt{\frac23}=\frac{4\sqrt6}{3}=K
\]
所以竖直直线也满足题设条件.

因此存在这样的直线 \(m\)，其方程为
\(x=-2\)，\(y=\frac{\sqrt3}{3}(x+2)\)，\(y=-\frac{\sqrt3}{3}(x+2)\).
\end{enumerate}
\end{solution}
